\section{Expert Annotation Details}
\label{sec:expert_annotation_details}
\subsection{Expert Annotation for Visual Context Questions}
We employed three computer science graduate students for annotating 300 visual context questions. Being provided with the full list of EMNLP 2023 papers, they were required to: (1) check that each anchor paper has arXiv documentation; (2) find figures or tables that contain comparative information with potential reasoning types described in Table \ref{tab:locality_reasoning}; (3) find the potential reference paper in the figure or table and ensure that it has arXiv documentation; and (4) write the visual context question. When they choose a figure or a table, they were required to fill in the corresponding visual context reasoning type as well as the ``direct answer'' to the visual context question. 

\subsection{Bias Mitigation for Visual Context Questions Annotation}
In preparation for the main annotation process, we conduct a pilot annotation stage where 20 papers where sampled. 
Annotators are instructed to generate three distinct questions per paper.  
These questions are subsequently analyzed by the authors and categorized into four distinct reasoning types: \emph{comparison}, \emph{data extraction}, \emph{location}, \emph{and visual understanding}.
These categories are comprehensive for scientific image understanding. 
By following the predefined reasoning type definitions in Table \ref{tab:locality_reasoning}, we mitigate the risk of annotator bias driven by their own preferences. 
Additionally, these reasoning types are not specific to NLP and are carefully chosen such that they are applicable in analyzing scientific images in the broader scientific fields. 

\subsection{Expert Annotation for Reference-based Questions}
We require each reference paper to have arXiv documentation. Then, we use the arXiv downloader to obtain the full text of the reference paper and generate subsequent reference-based questions (along with answers, explanations, and evidence) using the prompts described in Section \ref{prompt_detail}. We test these questions in the oracle setting, use GPT-based evaluators to evaluate if the answer generated in the oracle setting matches the answer generated along with the question. If they do not match, expert annotators proceeded to manually examine these questions and re-write the answers.

\subsection{Expert Annotation for Reference-based Reasoning}
In Section \ref{prompt_detail}, we automatically assign reasoning types concurrently with the generation of reference-based questions. To ensure the quality of the generated questions, we prompt GPT-4 with the question and its assigned reasoning type to ask if the question matches the reasoning type. For every question that GPT-4 flags as not matching the assigned reasoning type, expert annotators were instructed to manually examine the reasoning types and correct them when necessary. 

\subsection{Expert Annotation for Reference-based Answers}
\label{2_round_checking}

Following the two-round answer generation process mentioned in Section \ref{sec:benchmark_construction}, we manually checked 100 questions for which the first and second round answers matched in order to ensure the gold answers were indeed correct. Out of the 100 sampled questions, 96 questions were marked as correct by expert annotators, demonstrating the high-quality of \data benchmark. 
